\chapter*{Gedragscode en transparantieverklaring voor het gebruik van GenAI door KU Leuven-studenten (academiejaar 2025-2026)}
\addcontentsline{toc}{chapter}{Code of Conduct}
\markboth{Gedragscode}{Gedragscode}
Generatieve AI (GenAI)-tools kunnen worden gebruikt om verschillende soorten inhoud te genereren, zoals tekst, afbeeldingen, code, video, muziek of combinaties daarvan. Voorbeelden van dergelijke tools zijn onder andere ChatGPT, Google Gemini, Microsoft Copilot, Midjourney, Claude.ai, Perplexity.ai en DALL·E.
\newline
\newline
Deze gedragscode helpt studenten transparant te zijn over het gebruik van GenAI en sluit aan bij de universitaire principes van academische integriteit.

\subsection*{Belangrijke richtlijnen en opmerkingen}

\begin{itemize}
    \item \textbf{Gevoelige of persoonsgegevens:} Niet alle GenAI-tools beschermen je gegevens even goed. \textbf{Gebruik geen gratis GenAI-tools om vertrouwelijke of persoonsgegevens in te voeren.} Wat vertrouwelijke en persoonsgegevens zijn, vind je nader omschreven \href{https://www.kuleuven.be/onderwijs/learninglab/ondersteuning/genai/richtlijnen-veilig-gebruik-ai-tools#vertrouwelijkheidsniveaus}{in de drie vertrouwelijkheidsniveaus van het KU Leuven data classificatiemodel }(niet vertrouwelijk, vertrouwelijk, strikt vertrouwelijk). In geval van vertrouwelijke gegevens mag je enkel veilige tools gebruiken. Gebruik bij voorkeur \textbf{Microsoft 365 Copilot Chat met je KU Leuven account}; indien je toch een andere GenAI-tool wenst te gebruiken, moet je er \textbf{absoluut zeker} van zijn dat deze tool de informatie die je ingeeft \textbf{niet blijvend opslaat en hergebruikt. Vermijd het invoeren van deze gegevens;} doe het alleen voor zover strikt nodig en gebruik voor persoonsgegevens anonieme of gepseudonimiseerde gegevens. Voor strikt vertrouwelijke gegevens moet je bovendien eerst overleggen met je docent of promotor.
    
    \item \textbf{Auteursrechten:} Voor auteursrechtelijk beschermd materiaal (dat je rechtmatig bekomen hebt) gebruik je bij voorkeur \textbf{Microsoft 365 Copilot Chat met je KU Leuven account;} indien je toch een andere GenAI-tool wenst te gebruiken, moet je er \textbf{absoluut zeker} van zijn dat deze tool de informatie die je ingeeft \textbf{niet blijvend opslaat en hergebruikt.}
    
    \item GenAI mag niet gebruikt worden voor data of onderwerpen die onder een \textbf{geheimhoudingsverklaring} vallen (\href{https://www.kuleuven.be/stuvo/juridischedienst/studiestage/onderzoek}{Non-Disclosure agreement, NDA}). Zelfs indien er geen NDA is, kan het nodig zijn bepaalde informatie vertrouwelijk te behandelen, bijv. omwille van regelgevende vereisten of door het risico op ernstige schade voor de universiteit bij bekendmaking.  Twijfel je? Vraag je docent of promotor om advies.
    
    \item \href{https://iiw.kuleuven.be/studeren/masterproef/embargo}{Embargo op masterproef}: Overleg eerst met je docent of promotor of gebruik van GenAI is toegestaan.
    
    \item Overweeg steeds of het gebruik van een GenAI-tool verantwoord is, ook rekening houdend met \href{https://www.kuleuven.be/genai/duurzaamheid}{duurzaamheidsaspecten} (in het bijzonder bij gebruik als zoekmachine, taalhulp,…).
    
    \item \textbf Gebruik GenAI met een \textbf{wetenschappelijke en kritische houding.} De output is niet altijd correct.
    
    \item Je bent als student verantwoordelijk voor het naleven van artikel 84 van het Onderwijs- en Examenreglement. Je werk moet jouw eigen kennis, begrip en vaardigheden weerspiegelen. Plagiaatregels zijn ook van toepassing op werk dat (deels) met GenAI is gemaakt. 
    \newline
    \textbf{Artikel 84 OER:} \textit{“Elk gedrag van individuele studenten waardoor zij een juist oordeel omtrent de kennis, het inzicht en/of de vaardigheden van henzelf of van andere studenten geheel of gedeeltelijk onmogelijk maken of onmogelijk pogen te maken, wordt als een onregelmatigheid beschouwd die aanleiding kan geven tot een aangepaste sanctie. Een bijzondere vorm van een dergelijke onregelmatigheid is plagiaat, namelijk de overname zonder adequate bronvermelding van het werk (ideeën, teksten, structuren, ontwerpen, beelden, plannen, codes, ...) van anderen of van eerder werk van zichzelf, op identieke wijze of onder licht gewijzigde vorm. Elk bezit tijdens een examen van een niet-toegestaan hulpmiddel (zie artikel 65) wordt als een onregelmatigheid beschouwd.”}
    
    \item Om de academische integriteit te waarborgen en plagiaat te vermijden, vind je op deze website voor studenten \textbf{meer info over transparantie en correcte bronvermelding} bij gebruik van GenAI: \href{https://www.kuleuven.be/onderwijs/student/onderwijstools/artificiele-intelligentie#Transparantie}{NL}/\href{https://www.kuleuven.be/english/education/student/educational-tools/generative-artificial-intelligence#Transparency}{EN}.
    \item \textbf{Lees ook: KU Leuven richtlijnen voor verantwoord gebruik van GenAI tools, en bijkomende info }(\href{https://www.kuleuven.be/genai}{NL}/\href{https://www.kuleuven.be/english/genai}{EN})
\end{itemize}

\subsection*{Tot slot}
\begin{itemize}
    \item Twijfel je of je je gebruik van GenAI moet melden? Bespreek dit dan met je docent of promotor. Vermelden is altijd beter, zelfs indien het strikt genomen niet vereist is. Let op: vermelden van GenAI-gebruik betekent niet automatisch dat het ook toegestaan is. De rechterkolom in de tabel hieronder geeft daarover meer details (gedragscode). 
    \item Aangezien AI-tools zich snel ontwikkelen, hebben we nog niet op alles een pasklaar antwoord. Blijf dus op de hoogte van nieuwe ontwikkelingen, wees open maar ook kritisch, communiceer met docenten, assistenten en medestudenten, wees zo transparant mogelijk, en leer samen onderweg.
\end{itemize}
\newpage
% --- Student info row ---
% --TO BE FILLED IN--
\noindent
\textbf{Naam student:} \dotfill \hfill \textbf{Studentennummer:} \dotfill

\vspace{1em}

% --- Grey box section title ---
\noindent
\colorbox{lightgray}{\parbox{\dimexpr\linewidth-2\fboxsep\relax}{\textbf{Geef aan met “X” of het om een opdracht of bachelor-/masterproef gaat:}}}

\vspace{1em}

O Deze verklaring hoort bij een \textbf{opdracht.}

\begin{tabularx}{\linewidth}{@{}lXlX@{}}
\textbf{Naam opleidingsonderdeel:} & \dotfill & \textbf{Vakcode:} & \dotfill \\
\end{tabularx}

\vspace{1em}

Deze verklaring hoort bij mijn O \textbf{bachelorproef} of O \textbf{masterproef.}

\begin{tabularx}{\linewidth}{@{}lXlX@{}}
\textbf{Titel bachelor- of masterproef:} & \dotfill & \textbf{Promotor:} & \dotfill \\
\end{tabularx}

\vspace{1em}

\textbf{Dagelijkse begeleider:} \dotfill

\vspace{2em}

% --- Grey box section title ---
\noindent
\colorbox{lightgray}{\parbox{\dimexpr\linewidth-2\fboxsep\relax}{\textbf{Geef aan met "X":}}}

\vspace{1em}

O \textbf{Ik heb geen GenAI-tool gebruikt.} \\[1em]

O \textbf{Ik heb een GenAI-tool gebruikt.} In dit geval geef aan \textbf{welke tools}(bijv ChatGPT, M365 Copilot...): \dotfill

\newpage
% ====== PREAMBLE AANVULLINGEN ======

\definecolor{lightgray}{gray}{0.9}

% Kolomdefinitie: mooie linkslijnende X-kolom
\newcolumntype{Y}{>{\raggedright\arraybackslash}X}

% Veilige cel-wrapper (geen \setlist hier!)
\newcommand{\Cell}[1]{\begin{minipage}[t]{\linewidth}\raggedright\small #1\end{minipage}}

% (optioneel) iets compactere tabel
\setlength{\tabcolsep}{6pt}
\renewcommand{\arraystretch}{1.2}

\begin{longtable}{|p{0.25\textwidth}|p{0.2\textwidth}|p{0.55\textwidth}|}
\hline 
\textbf{GenAI-tool gebruikt als/voor:} & \textbf{Naam van de gebruikte GenAI-tool. } Indien nuttig, omschrijf hoe je GenAI gebruikte, in verhouding tot de gedragscode. & \textbf{Gedragscode: \newline
Eerbiedig voor elk van onderstaande categorieën altijd de hoger genoemde belangrijke richtlijnen en opmerkingen (bijv. auteursrechtelijk beschermd material, gevoelige of persoonsgegevens,…)} \\ \hline
\endfirsthead

%\hline
%\textbf{GenAI assistance used as/for:} & \textbf{Name of the GenAI tool(s) used} If helpful, also describe in which way you were using GenAI related to what is specified as code of conduct. & \textbf{Code of conduct:} For each of the categories below, always take into account the important guidelines and remarks mentioned above (e.g. copyrighted data, sensitive or personal data,…). \\  \\ \hline
%\endhead

\hline
\endfoot

\hline
\endlastfoot

\textbf{Als taalhulp voor het nakijken of verbeteren van zelf geschreven teksten} & \textit{voeg hier tekst toe} & Dit gebruik is gelijkaardig aan tools voor spelling- of grammaticacontrole. In het algemeen hoef je in de tekst niet te verwijzen naar een dergelijk GenAI-gebruik.  

\textbf{Echter, let op:  }
\begin{itemize}
    \item Indien je de GenAI-tool gebruikt om teksten te verbeteren die je niet zelf schreef, \textbf{moet} je \href{https://bib.kuleuven.be/training-en-tutorials/citeren/basisprincipes}{verwijzen} naar \textbf{de originele bron of auteur} (indien niet, pleeg je \href{https://www.kuleuven.be/onderwijs/plagiaat/index}{plagiaat} en dus een onregelmatigheid).
\end{itemize} \\ \hline
\textbf{Als parafraseertool} & \textit{voeg hier tekst toe} & Dit gebruik is toegestaan, tenzij het is verboden door de docent of de opleiding. Je mag je eigen teksten of die van andere auteurs parafraseren en je laten inspireren door wat de GenAI-tool voorstelt (tenzij het is verboden). In het algemeen hoef je in de tekst niet te verwijzen naar een dergelijk GenAI-gebruik.

\textbf{Echter, let op:} 
\begin{itemize}
    \item Indien je de GenAI-tool gebruikt om teksten te parafraseren die je niet zelf schreef, \textbf{moet} je \href{https://bib.kuleuven.be/training-en-tutorials/citeren/basisprincipes}{verwijzen} naar \textbf{de originele bron of auteur }(indien niet, pleeg je \href{https://www.kuleuven.be/onderwijs/plagiaat/index}{plagiaat} en dus een onregelmatigheid).
    \item \textbf{Kopieer geen hele tekstblokken}; je moet immers begrijpen en kritisch verwerken wat je schrijft.
    \item Controleer altijd de correctheid en betekenis van de tekst.
\end{itemize} \\ \hline
\textbf{Als vertaalhulp om teksten te verbeteren die je zelf schreef of om teksten van anderen beter te begrijpen} & \textit{voeg hier tekst toe} & Dit gebruik is toegestaan, tenzij het is verboden door de docent of de opleiding. Het is gelijkaardig aan het gebruik van vertaaltools (Google translate, DeepL,…).

In het algemeen hoef je in de tekst niet te verwijzen naar een dergelijk GenAI-gebruik.
\textbf{Echter, let op: } 
\begin{itemize}
    \item Indien je de GenAI-tool gebruikt om teksten te vertalen die je niet zelf schreef, \textbf{moet} je \href{https://bib.kuleuven.be/training-en-tutorials/citeren/basisprincipes}{verwijzen} naar \textbf{de originele bron of auteur} (indien niet, pleeg je \href{https://www.kuleuven.be/onderwijs/plagiaat/index}{plagiaat} en dus een onregelmatigheid).
    \item Controleer altijd de correctheid en betekenis van de vertaalde tekst.
\end{itemize} \\ \hline
\textbf{Als een zoekmachine om informatie te vinden over een topic of om hierover bestaand onderzoek te vinden} & \textit{voeg hier tekst toe} & Dit gebruik is gelijkaardig aan bijv. Google search of Wikipedia. Als je je eigen tekst schrijft op basis van deze info, hoef je niet te verwijzen naar een dergelijk GenAI-gebruik. 
\textbf{Echter, let op: }
\begin{itemize}
    \item Je \textbf{moet} wel \href{https://bib.kuleuven.be/training-en-tutorials/citeren/basisprincipes}{verwijzen} naar het \textbf{bestaande onderzoek en de verwijzingen} die je hebt gecheckt en gebruikt (indien niet, pleeg je \href{https://www.kuleuven.be/onderwijs/plagiaat/index}{plagiaat} en dus een onregelmatigheid).
    \item -	Wees je ervan bewust dat de output van de GenAI-tool niet 100\% betrouwbaar is. De output kan niet helemaal correct of beperkt zijn door de database die de tool gebruikt. Bovendien evolueert kennis; mogelijk is de database van de GenAI-tool niet up-to-date. Controleer daarom altijd de informatie en kopieer het niet gewoon; je moet immers begrijpen en kritisch verwerken wat je schrijft.
\end{itemize} \\ \hline
\textbf{Voor literatuuronderzoek} & \textit{voeg hier tekst toe} & Dit gebruik is gelijkaardig aan bijv. Google Scholar. In het algemeen hoef je in de tekst niet te verwijzen naar een dergelijk GenAI-gebruik.

\medskip
\textbf{Echter, let op: }
\begin{itemize}
    \item Je \textbf{moet} wel \href{https://bib.kuleuven.be/training-en-tutorials/citeren/basisprincipes}{verwijzen} naar de \textbf{literatuur verwijzingen} die je hebt gecheckt en gebruikt (indien niet, pleeg je \href{https://www.kuleuven.be/onderwijs/plagiaat/index}{plagiaat} en dus een onregelmatigheid).
    \item Wees je ervan bewust dat het zoekresultaat beperkt is tot de database van de GenAI-tool. Zoek en analyseer daarom vervolgens de \textbf{originele wetenschappelijke bronnen}. Verifieer, interpreteer, analyseer en verwerk de informatie die je bekwam. Dus niet louter copy-pasten; je moet immers begrijpen en kritisch analyseren wat je schrijft.
    \item Wees je ervan bewust dat sommige GenAI-tools \href{https://www.kuleuven.be/onderwijs/learninglab/ondersteuning/genai/richtlijnen-veilig-gebruik-ai-tools}{geen of verkeerde verwijzingen} aanleveren. Als student ben je verantwoordelijk voor het checken van het ontbreken en de correctheid van de verwijzingen. Opnieuw: dus niet louter copy-pasten.
\end{itemize} \\ \hline
\textbf{Voor het genereren van programmeercode} & \textit{voeg hier tekst toe} & GenAI-gebruik voor het genereren van code is toegestaan, tenzij het is verboden door de docent of opleiding. Vermeld altijd het gebruik van de GenAI-tool, in overeenstemming met de richtlijnen \href{https://www.kuleuven.be/onderwijs/student/onderwijstools/artificiele-intelligentie#Transparantie}{op de pagina hoe transparant te zijn over het gebruik van GenAI.}
\\ \hline
\textbf{Voor het genereren van nieuwe (onderzoeks)ideeën} & \textit{voeg hier tekst toe} & Dit GenAI-gebruik is toegestaan, tenzij het is verboden door de docent of opleiding. Vermeld altijd het gebruik van de GenAI-tool, in overeenstemming met de richtlijnen op \href{https://www.kuleuven.be/onderwijs/student/onderwijstools/artificiele-intelligentie#Transparantie}{op de pagina hoe transparant te zijn over het gebruik van GenAI.}

\medskip
\textbf{Echter, let op:}
\begin{itemize}
    \item Ga altijd na of het idee effectief nieuw is of niet. Er is een grote kans dat het gebaseerd is op bestaand werk. Indien dat het geval is, \textbf{moet} je \href{https://bib.kuleuven.be/training-en-tutorials/citeren/basisprincipes}{verwijzen} naar het \textbf{bestaande werk }(indien niet, pleeg je \href{https://www.kuleuven.be/onderwijs/plagiaat/index}{plagiaat} en dus een onregelmatigheid).
\end{itemize}\\ \hline
\textbf{Voor het genereren van synthetische data } & \textit{voeg hier tekst toe} & GenAI-gebruik voor het genereren van synthetische data is toegestaan, voor zover methodologisch en ethisch verantwoord. tenzij het is verboden door de docent of opleiding. Vermeld altijd het gebruik van de GenAI-tool, in overeenstemming met de richtlijnen \href{https://www.kuleuven.be/onderwijs/student/onderwijstools/artificiele-intelligentie#Transparantie}{op de pagina hoe transparant te zijn over het gebruik van GenAI.}

\medskip
\textbf{Echter, let op:}
\begin{itemize}
    \item Check altijd de gegenereerde data op kwaliteit en mogelijke bias. De kwaliteit van de output is immers sterk afhankelijk van de kwaliteit van de data waarop de modellen getraind zijn.
\end{itemize} \\ \hline
\textbf{Voor het genereren van tekstblokken (anders dan het hoger toegelaten gebruik zonder verwijzing)} & \textit{voeg hier tekst toe} & Volgens artikel 84 Onderwijs- en Examenreglement moet je tekst een juist oordeel toelaten over je eigen kennis, inzicht en vaardigheden. Daarom is het \textbf{verboden} tekstblokken op te nemen zonder aanhalingstekens en verwijzing naar het gebruik van een GenAI-tool.

\medskip
\textbf{Echter, let op:}
\begin{itemize}
    \item Indien het echt nodig is tekstblokken van een GenAI-tool op te nemen, bijv. omwille van de aard van je opdracht, \textbf{citeer het door gebruik van aanhalingstekens en vermeld het gebruik van de GenAI-tool,} in overeenstemming met de richtlijnen \href{https://www.kuleuven.be/onderwijs/student/onderwijstools/artificiele-intelligentie#Transparantie}{op de pagina hoe transparant te zijn over het gebruik van GenAI.}
    \item In het algemeen moet een dergelijk GenAI-gebruik tot een \textbf{absoluut minimum} beperkt worden; check ook altijd de originele bronnen.
\end{itemize} \\ \hline
\textbf{Voor het genereren van beeldmateriaal, video of audio} & \textit{voeg hier tekst toe} & Gebruik van GenAI voor het genereren van afbeeldingen, video of audio is toegestaan, tenzij het verboden is door de docent of opleiding. 

\medskip
\textbf{Echter, let op:}
\begin{itemize}
    \item \textbf{Vermeld het gebruik van de GenAI-tool,} in overeenstemming met de richtlijnen op \href{https://www.kuleuven.be/onderwijs/student/onderwijstools/artificiele-intelligentie#Transparantie}{de pagina hoe transparant te zijn over het gebruik van GenAI.}
    \item -	Indien je met de GenAI-tool werkt op basis van bestaande afbeeldingen, video of audio, moet je \href{https://bib.kuleuven.be/training-en-tutorials/citeren/basisprincipes}{verwijzen} naar het \textbf{bestaande materiaal} (indien niet, pleeg je \href{https://www.kuleuven.be/onderwijs/plagiaat/index}{plagiaat} en dus een onregelmatigheid). Leg het GenAI-gebruik uit in de methodesectie (als die er is) en vermeldt eventueel ook de gebruikte prompts (of een link ernaar) en de volledige GenAI-output (of de geschiedenis). 
\end{itemize} \\ \hline
\textbf{Overig gebruik (leg uit; het kan ook gaan om een combinatie van bovenstaande vormen)} & \textit{voeg hier tekst toe} & Neem \textbf{vooraf} contact op met het onderwijsteam van het opleidingsonderdeel of de promotor van je thesis, om te checken of je voorgenomen GenAI-gebruik is toegestaan en welke voorwaarden eventueel van toepassing zijn. Licht ook de programmadirecteur in.

Motiveer hoe het voorgenomen gebruik in overeenstemming is met de opdracht of thesis en met art. 84 Onderwijs- en Examenreglement en \href{https://www.kuleuven.be/genai}{de KU Leuven richtlijnen voor verantwoord gebruik van GenAI tools.} Leg ook de meerwaarde van het GenAI-gebruik uit. 

Naargelang het GenAI-gebruik, kan het nodig zijn ernaar te verwijzen in je tekst, volgens de de richtlijnen \href{https://www.kuleuven.be/onderwijs/student/onderwijstools/artificiele-intelligentie#Transparantie}{op de pagina hoe transparant te zijn over het gebruik van GenAI}
\\ \hline
\end{longtable}